\subsection{Cognitive runtime and admission contract}
The model adapter returns a structured reply containing public text, tool calls, terminal reason and per-round metrics. This boundary permits scripted providers in integration tests while keeping the production Ollama transport responsible for streamed protocol validation. It does not make every provider interchangeable: the admitted request remains bound to a specific profile and model deployment. The adapter parses newline-delimited JSON, validates terminal metadata, limits public output, and does not copy the provider's thinking field into user-facing text or diagnostics. Streamed text is provisional until a valid terminal reply arrives; a late protocol failure cannot retract text already displayed.

ContextBuilder clones caller-owned JSON and validates message/tool structure before rendering. It associates messages with whole user-turn groups, so eviction removes old groups rather than separating a tool result from the interaction that produced it. Evidence is explicitly marked as untrusted data in the system payload. For rendered token count $T$, profile safety allowance $S$, output reserve $O$, and configured context $C$, admission requires
\begin{equation}
T + S + O \leq C.
\end{equation}
Here $T$ includes template rendering and the supplied tool schemas, not merely the lengths of visible message strings. The resulting immutable snapshot carries a fingerprint, profile identity, deployment revisions, token counts and evicted message identifiers. Send-time validation checks that these have not changed. Such consistency checks are distinct from authority checks: a correctly fingerprinted request does not grant permission to execute its proposed actions.

When mandatory system/evidence/tool content and the newest group cannot fit, the builder returns context\_overflow. Persist-mode admission additionally tries bounded history reductions and an extractive continuity note. Working notes are reinserted from turn state each round, so trimming old tool output need not remove every recorded finding. These are lossy context-management strategies, not lossless history compression. Provider prompt counts are inspected at the terminal boundary: a count above the admitted value is rejected, while a smaller count is logged rather than rejected. Thus the narrow historical exact-agreement calibration must not be generalized into an unconditional guarantee of equality for every production request.

\subsection{Dispatch, approvals and structured outcomes}
ToolRegistry checks the current execution context and capability, snapshots arguments before awaiting, and performs strict typed validation. Registered tools carry an argument model, capability, risk and asynchronous handler. The control store prepares a durable operation; destructive actions can enter approval\_pending according to configured exemptions, while arbitrary-risk operations require the appropriate execution authority. Approval therefore concerns the registered operation and arguments rather than a natural-language statement that the model has obtained permission.

An EffectGuard connects handler work to current control state, operation identity and deadline. Handlers are trusted cooperative code and must recheck before additional effects. This arrangement limits stale authority within the managed path; it cannot contain malicious in-process code or replace an operating-system sandbox. Handler outcomes retain status, error code, bounded content, truncation and side-effect information. A cancellation or timeout after an operation starts can leave an unknown outcome: stopping the coroutine does not prove that a remote request or filesystem effect never happened. Durable no-replay state and explicit inspection are therefore preferable to an automatic retry that could duplicate an effect.

The agent consumes these structured results across bounded model rounds. Its terminal result distinguishes completion, failure, cancellation and an approval wait, and includes persistence state, evidence identifiers, events and model metrics. Interrupted work may retain an explicitly unfinished transcript; it must not be relabeled successful because partial text or a file exists. Selected integration tests exercise these contracts using controlled handlers and providers, while natural model reasoning remains a separate empirical question.

